\documentclass[10pt,a4paper]{amsart}
\usepackage[margin=19mm]{geometry}
\usepackage{amsmath,amssymb,mathtools}
\usepackage[scr=rsfs,cal=euler]{mathalfa}
\usepackage{xfrac}
\usepackage[unicode,hidelinks,bookmarks=false]{hyperref}
\newcommand{\ie}{i.e.\ }
\newcommand{\cf}{cf.\ }
\newcommand{\resp}{respectively\ }
\newcommand{\Subsection}{\subsection}
\providecommand{\nobreakdash}{\nobreak\hbox{-}}
\setlength{\emergencystretch}{2em}
\multlinegap0pt
\usepackage{bm}
\usepackage{bbm}
\usepackage{dsfont}
\usepackage{xspace}
\usepackage{cancel}
\usepackage{mathtools}
\usepackage{amsthm}
\makeatletter
\@ifundefined{lem}{\newtheorem{lem}{Lemma}}{}
\makeatother
\DeclareMathOperator{\nrd}{nrd}
\DeclareMathOperator{\trd}{trd}
\newcommand{\D}{\mathcal{D}}
\newcommand{\J}{\mathfrak{J}}
\newcommand{\M}{\mathfrak{M}}
\newcommand{\p}{\mathfrak{p}}
\renewcommand{\O}{\mathcal{O}}
\renewcommand{\P}{\mathfrak{P}}
\renewcommand{\o}{\mathfrak{o}}
\title[Chebotarev geodesic theorem: non-split case]{Chebotarev geodesic theorem: non-split case}
\author{Alberto Acosta Reche}
\date{Source: arXiv:2608.04923v1 (CC BY 4.0)}
\begin{document}
\maketitle

\noindent\emph{Source and licence.} Chebotarev geodesic theorem: non-split case. Version arXiv:2608.04923v1 (CC BY 4.0), distributed under the Creative Commons Attribution 4.0 International licence, as recorded in the arXiv licence metadata for this exact version and shown on the abstract page https://arxiv.org/abs/2608.04923v1. Full rights notice: licenses/arxiv-2608.04923-CC-BY-4.0.txt.\par\smallskip

\noindent\textit{Editorial scope.} Counted unit: the Lem whose source label is \texttt{lem:charpolyramified} in the article (source0.tex lines 813--820, source label \texttt{lem:charpolyramified}), delivered together with the article's own complete original proof of it (source0.tex lines 821--861, one \texttt{proof} environment of 41 lines). No mathematics is written, repaired or re-derived: statement and proof are the article's own text line for line. Required context retained uncounted: eq:tracecyclicalgebra (lines 650--653); eq:normcyclicalgebra (lines 655--657); eq:descriptionordercyclicpart1 (lines 659--661); auxeq:descriptionuniformizers (lines 707--709); lem:imageofhhereditaryorders (lines 758--760); lem:preimagesramifiedunramified (lines 765--767); lem:auxiliaryhensel (lines 773--775). Every definition, display and label of those lines is retained unchanged. Omitted from this package and not redistributed: 1--649, 654--654, 658--658, 662--706, 710--757, 761--764, 768--772, 776--812, 862--1588. Nothing inside a retained excerpt is omitted or altered: every retained body is delivered exactly as the article's lines are written, so no omission receipt of any kind accompanies this package. Citations: the retained excerpts carry no citation command at all, so no bibliography entry of the article is transcribed and no reference is redistributed. Reference adaptation: none was needed and none was made; every reference inside the delivered unit resolves inside it, so no unresolved reference or citation is produced. Printed slips kept literal: no printed word, symbol, hypothesis or number of the counted statement or of its proof was altered. Layout and class: the article's own class is \texttt{article} with packages including amsmath, amsthm, amssymb, mathtools, leftindex, tikz-cd, verbatim, bbm, enumitem, parskip, geometry, biblatex, hyperref; the wrapper is the lane's pinned \texttt{amsart} editorial wrapper at 10pt on A4, which restates the theorem-like environments the retained text needs and adds the article's own macro definitions that the retained text needs (\texttt{\textbackslash D}, \texttt{\textbackslash J}, \texttt{\textbackslash M}, \texttt{\textbackslash O}, \texttt{\textbackslash P}, \texttt{\textbackslash nrd}, \texttt{\textbackslash o}, \texttt{\textbackslash p}, \texttt{\textbackslash trd}), with extra wrapper packages (bm, bbm, dsfont, xspace, cancel, mathtools). The delivered numbering is the unit's own. Assets: no retained span contains a figure environment, a graphics-inclusion command, a PDF-page inclusion, a table or a TikZ picture; no image file is copied into this package, the package declares no asset dependency and no image file is redistributed with it. Licence: version arXiv:2608.04923v1 of this article is distributed under the Creative Commons Attribution 4.0 International licence, as recorded in the arXiv licence metadata for this exact version and shown on the abstract page https://arxiv.org/abs/2608.04923v1. A full rights notice is delivered at licenses/arxiv-2608.04923-CC-BY-4.0.txt. Characters: every retained excerpt is pure ASCII, so the delivered engine, XeLaTeX with its default Latin Modern text font, reports no missing character.
\begin{equation}\label{eq:tracecyclicalgebra}
    \trd(x + \varpi_\O y) = \trd(x),\quad 
    \overline{x + \varpi_\O y} = \overline{x} - \varpi_\O y,
\end{equation}

\begin{equation}\label{eq:normcyclicalgebra}
    \nrd(x + \varpi_\O y) = \nrd(x) -\varpi \nrd(y),
\end{equation}

\begin{equation}\label{eq:descriptionordercyclicpart1}
    \O = \o_{E_B} \oplus \varpi_\O \o_{E_B}, \quad \O^\times = \o_{E_B}^\times \oplus \varpi_\O \o_{E_B},
\end{equation}

\begin{equation}\label{auxeq:descriptionuniformizers}
\varpi_\O \O^\times = \varpi \o_E \oplus \varpi_\O \o_{E}^\times.
\end{equation} 

\begin{lem}\label{lem:imageofhhereditaryorders}
We have $Z_u \cup Z_r \subset h(\M)\cap h(\O_\D)$, as well as $Z_r \subset h(\J)$.
\end{lem}

\begin{lem}\label{lem:preimagesramifiedunramified}
If $\O \in \{\J, \O_\D\}$ we have $h_\O^{-1}(Z_r) = \varpi_\O \O^\times$. If $\O = \O_\D$, we have $h_{\O_\D}^{-1}(Z_u) = U_\D - \o^\times U_\D^1$. 
\end{lem}

\begin{lem}\label{lem:auxiliaryhensel}
Let $f: \o^n \rightarrow \o^m$ be a map of the form $f(x) = Ax + \varpi P(x)$, where $A \in M_{m \times n}(\o)$ and $P: \o^n \rightarrow \o^m$ is a polynomial map with coefficients in $\o$. Then $f$ is surjective if and only if $x\mapsto Ax$ is surjective modulo $\p$.
\end{lem}

\begin{lem}\label{lem:charpolyramified}
Let $\O \in \{\J, \O_\D\}$ and $\alpha \in h_\O^{-1}(Z_r)$. Then, we have: 
\begin{enumerate}
    \item [i)]$h(\alpha U_\O) = Z_r = \p \times (\p \smallsetminus \p^2)$,
    \item [ii)]$h(\alpha U_\O^{2n - 1}) = h(\alpha) + \p^n \times \p^{n + 1}$ for $n \geq 1$, and
    \item [iii)]$h(\alpha U_\O^{2n}) = h(\alpha) + \p^{n + 1}\times \p^{n + 1}$ for $n \geq 1$.
\end{enumerate}
\end{lem}

\begin{proof}
Recall that $U_\O = \O^\times$. By Lemma \ref{lem:preimagesramifiedunramified} we know that $h_\O^{-1}(Z_r) = \varpi_\O U_\O$. Since $\alpha \in h_\O^{-1}(Z_r)$ by hypothesis, we deduce that $h_\O^{-1}(Z_r) = \alpha U_\O$. On the other hand, by Lemma \ref{lem:imageofhhereditaryorders} we know that $Z_r \subset h(\O)$. Part i) follows from these observations. 

We now deal with part ii). Note that $\alpha U_\O^{2n-1} = \alpha + \P_\O^{2n} = \alpha + \varpi^n \O$. Write $\beta \in \alpha U_\O^{2n-1}$ as $\beta = \alpha + \varpi^n y$, with $y \in \O$. Furthermore, recalling \eqref{eq:descriptionordercyclicpart1} and \eqref{auxeq:descriptionuniformizers} we write
\begin{equation}\label{auxeq:decompositionelements}
    \alpha = \varpi \gamma + \varpi_\O \delta, \quad y = z + \varpi_\O w
\end{equation}
with $\gamma \in \o_E, \delta \in \o_E^\times$ and $z, w \in \o_E$. Recall that $E = F\times F$ if $\O = \J$, and $E/F$ is the unramified quadratic extension if $\O = \O_\D$. Using \eqref{eq:tracecyclicalgebra} and \eqref{eq:normcyclicalgebra} we see that 
\begin{equation*}
    \trd(y) = \trd(z), \quad \trd(\overline{\alpha}y) = \varpi \trd(\overline{\gamma}z) - \varpi \trd(\overline{\delta}w),\quad \nrd(y) = \nrd(z)-\varpi \nrd(w).
\end{equation*}
Therefore, 
\begin{equation*}
\begin{aligned}
\trd(\beta) & = \trd(\alpha) + \varpi^n \trd(z),\\
\nrd(\beta) & = \nrd(\alpha) + \varpi^{n + 1} (\trd(\overline{\gamma}z)- \trd(\overline{\delta}w)) + \varpi^{2n}\nrd(z) - \varpi^{2n + 1}\nrd(w).
\end{aligned}
\end{equation*}
From this equation we see that $h(\alpha U_\O^{2n-1})\subset h(\alpha) + \p^n \times \p^{n + 1}$. To show equality in this inclusion, we need to prove that the map $g: \o_E^2 \rightarrow \o^2$ given by
\begin{equation*}
    g(z, w) = (\trd(z), \trd(\overline{\gamma}z - \overline{\delta}w) + \varpi^{n-1}\nrd(z) - \varpi^{n}\nrd(w))
\end{equation*}
is surjective. Let $(t, n) \in \o^2$. Since $\trd(\o_E) = \o$, we can always choose $z$ so that $\trd(z) = t$. After $z$ is fixed, the task reduces to showing that the map $g': \o_E \rightarrow \o$ given by
\begin{equation*}
    g'(w) = \trd(\overline{\delta}w) + \varpi^n \nrd(w)
\end{equation*}
is surjective. Since $\delta \in \o_E^\times$, we have $\trd(\overline{\delta}\o_E) = \o$, so the surjectivity of $g'$ follows from Lemma \ref{lem:auxiliaryhensel}. This finishes the proof of part ii).

For part iii), we have $\alpha U_\O^{2n} = \alpha + \P_\O^{2n + 1}$. If we write $\beta = \alpha + \varpi^n \varpi_\O y$, with $y \in \O$, then 
\begin{equation*}
    \trd(\beta) = \trd(\alpha) + \varpi^n \trd(\varpi_\O y), \quad \nrd(\beta) = \nrd(\alpha) + \varpi^n \trd(\overline{\alpha}\varpi_\O y) - \varpi^{2n + 1} \nrd(y).
\end{equation*}
Since $\trd(\P_\O) = \p$, it is clear that $h(\alpha U_\O^{2n}) \subset h(\alpha) + \p^{n + 1} \times \p^{n + 1}$. To check equality, we use the decomposition \eqref{auxeq:decompositionelements} and find that
\begin{equation*}
\begin{aligned}
    \trd(\beta) & = \trd(\alpha) + \varpi^{n + 1} \trd(w),\\
    \nrd(\beta) & = \nrd(\alpha) -\varpi^{n+1}\trd(\overline{\delta}z) + \varpi^{n + 2}\trd(\overline{\gamma}w) - \varpi^{2n + 1}(\nrd(z) - \varpi \nrd(w)).
\end{aligned}
\end{equation*} 
By Lemma \ref{lem:auxiliaryhensel}, it is now enough to show that the map $g: \o_E^2 \rightarrow \o^2$ given by $g(z, w) = (\trd(w), -\trd(\overline{\delta}z))$ is surjective. This is clear since $\delta \in \o_E^\times$ and $\trd(\o_E) = \o$.
\end{proof}

\end{document}
